% ================================================================
%  sesion02.tex  —  Sesión 2 de 20
%  ¿Qué se puede medir y representar en una ciudad?
%  Almería en Cifras y Formas · Matemáticas 1.º ESO
%  Compilar con:  pdflatex sesion02.tex  (dos pasadas)
% ================================================================
\documentclass[aspectratio=169, 10pt, spanish]{beamer}
\usepackage{almeria-beamer}

% ---------------------------------------------------------------
%  DATOS DE LA SESIÓN
% ---------------------------------------------------------------
\renewcommand{\SessionNumber}{02}
\renewcommand{\SessionTitle}{¿Qué se puede medir y representar en una ciudad?}
\renewcommand{\SessionName}{Sesión 02}
\renewcommand{\SessionObjective}{%
  Clasificar los tipos de información matemática presentes en una ciudad
  (ubicación, medida, dato) e identificar ejemplos reales en Almería.}
\renewcommand{\BlockName}{Bloque 1: Mapa y coordenadas}

% ================================================================
\begin{document}

% ---------------------------------------------------------------
%  PORTADA
% ---------------------------------------------------------------
\begin{frame}[plain]
  \titlepage
\end{frame}

% ---------------------------------------------------------------
%  ÍNDICE DE LA SESIÓN
% ---------------------------------------------------------------
\begin{frame}{Índice de la sesión}
  \begin{columns}[t]
    \begin{column}{0.5\textwidth}
      \begin{enumerate}
        \item Tres tipos de información matemática
        \item Magnitudes y unidades
        \item Tabla de conversión de unidades
        \item Ejemplos reales de Almería
        \item Conversiones con la Rambla y el Puerto
      \end{enumerate}
    \end{column}
    \begin{column}{0.5\textwidth}
      \begin{enumerate}
        \setcounter{enumi}{5}
        \item Ejercicios multinivel (Bloom)
        \item Evidencia del día
        \item Error frecuente
        \item Resumen
      \end{enumerate}
    \end{column}
  \end{columns}
\end{frame}

% ---------------------------------------------------------------
\seccionframe{1. Tres tipos de información matemática}
% ---------------------------------------------------------------

\begin{frame}{¿Qué información matemática hay en una ciudad?}
  \begin{columns}[c]
    \begin{column}{0.52\textwidth}
      \begin{teoriabox}
        Toda la información matemática de una ciudad
        puede clasificarse en \textbf{tres grandes tipos}:
        \begin{itemize}
          \item \textbf{Ubicación:} ¿Dónde está?
          \item \textbf{Medida:} ¿Cuánto mide/pesa/dura?
          \item \textbf{Dato:} ¿Cuántos/cuántas en un momento?
        \end{itemize}
      \end{teoriabox}
      \vspace{6pt}
      Reconocer el tipo correcto nos ayuda a elegir
      la \textbf{herramienta matemática} adecuada.
    \end{column}
    \begin{column}{0.44\textwidth}
      \begin{center}
        \begin{tikzpicture}[scale=0.88,
            nodo/.style={rounded corners=5pt, minimum width=2.5cm,
                         minimum height=0.7cm, align=center,
                         font=\AlmFontSmall\bfseries}]
          % Nodo raíz
          \node[nodo, fill=AlmAzul, text=AlmBlanco]
            (inf) at (0,2.8) {Información\\matemática};
          % Tres hijos
          \node[nodo, fill=AlmAzulM, text=AlmBlanco]
            (ubi) at (-2.6,1.2) {Ubicación};
          \node[nodo, fill=AlmAmbar!80!black, text=AlmBlanco]
            (med) at (0,1.2)    {Medida};
          \node[nodo, fill=AlmVerde, text=AlmBlanco]
            (dat) at (2.6,1.2)  {Dato};
          % Flechas
          \draw[AlmAmbar, line width=1.4pt, ->]
            (inf) -- (ubi);
          \draw[AlmAmbar, line width=1.4pt, ->]
            (inf) -- (med);
          \draw[AlmAmbar, line width=1.4pt, ->]
            (inf) -- (dat);
          % Ejemplos pequeños
          \node[font=\AlmFontTiny, color=AlmGris, align=center]
            at (-2.6,0.45) {coordenadas\\referencias};
          \node[font=\AlmFontTiny, color=AlmGris, align=center]
            at (0,0.45)    {longitud\\superficie\\masa};
          \node[font=\AlmFontTiny, color=AlmGris, align=center]
            at (2.6,0.45)  {población\\temperatura\\turistas};
        \end{tikzpicture}
      \end{center}
    \end{column}
  \end{columns}
\end{frame}

% -------
\begin{frame}{Los tres tipos en detalle}
  \begin{columns}[T]
    \begin{column}{0.32\textwidth}
      \begin{definicionbox}{Ubicación}
        \AlmFontFootnote
        Indica \textbf{dónde} se encuentra algo
        mediante coordenadas o referencias espaciales.\\[4pt]
        Ejemplo en Almería:\\
        \datoalm{Alcazaba}{36,8409°\,N\;/\;2,4699°\,O}
      \end{definicionbox}
    \end{column}
    \begin{column}{0.32\textwidth}
      \begin{definicionbox}{Medida}
        \AlmFontFootnote
        Expresa una \textbf{magnitud física} con un valor
        \emph{fijo} (no cambia con el tiempo).\\[4pt]
        Ejemplo en Almería:\\
        \datoalm{Rambla}{370 m de longitud}\\
        \datoalm{Puerto}{35 ha de superficie}
      \end{definicionbox}
    \end{column}
    \begin{column}{0.32\textwidth}
      \begin{definicionbox}{Dato}
        \AlmFontFootnote
        Información numérica \textbf{variable}: puede
        cambiar según el momento o el año.\\[4pt]
        Ejemplo en Almería:\\
        \datoalm{Temperatura}{máx. 38 °C en agosto}\\
        \datoalm{Turistas}{5,8 millones en 2023}
      \end{definicionbox}
    \end{column}
  \end{columns}
  \vspace{8pt}
  \begin{notabox}[Clave para distinguirlos]
    \AlmFontSmall
    \correcto\ La \textbf{medida} no cambia: la longitud de la Rambla
    siempre es 370 m.\\
    \correcto\ El \textbf{dato} cambia: la temperatura de mañana puede ser distinta.
  \end{notabox}
\end{frame}

% ---------------------------------------------------------------
\seccionframe{2. Magnitudes y unidades}
% ---------------------------------------------------------------

\begin{frame}{Magnitudes físicas y sus unidades}
  \begin{columns}[c]
    \begin{column}{0.52\textwidth}
      \begin{teoriabox}
        Una \textbf{magnitud} es todo aquello que se puede
        medir. Toda medida consta de:
        \[
          \text{medida} = \underbrace{\text{número}}_{\text{valor}}
            \;\times\;
          \underbrace{\text{unidad}}_{\text{patrón}}
        \]
        Ejemplo: $370\ \unidad{m}$ — el número es 370
        y la unidad es el metro (\unidad{m}).
      \end{teoriabox}
      \vspace{6pt}
      \begin{notabox}[Regla de oro]
        \AlmFontSmall
        Nunca escribas un resultado \textbf{sin unidades}.\\
        «370» no significa nada; «$370\ \text{m}$» sí.
      \end{notabox}
    \end{column}
    \begin{column}{0.44\textwidth}
      \vspace{4pt}
      {\AlmFontSmall
      \begin{tabular}{@{}lll@{}}
        \toprule
        \textbf{Magnitud} & \textbf{Unidad base} & \textbf{Símbolo}\\
        \midrule
        Longitud     & metro        & \unidad{m}\\
        Superficie   & metro cuadrado & \unidad{m\textsuperscript{2}}\\
        Masa         & kilogramo    & \unidad{kg}\\
        Tiempo       & segundo      & \unidad{s}\\
        Temperatura  & grado Celsius & \unidad{°C}\\
        \bottomrule
      \end{tabular}}
      \vspace{8pt}
      \begin{ejemplobox}
        \AlmFontFootnote
        La \textbf{Catedral de Almería} tiene
        55 \unidad{m} de altura. Su torre mide
        un \textbf{tiempo} de sombra diferente
        según la estación del año.
      \end{ejemplobox}
    \end{column}
  \end{columns}
\end{frame}

% ---------------------------------------------------------------
\seccionframe{3. Tabla de conversión de unidades}
% ---------------------------------------------------------------

\begin{frame}{Conversiones de unidades más usadas}
  \begin{columns}[c]
    \begin{column}{0.56\textwidth}
      \begin{center}
        {\AlmFontSmall
        \begin{tabular}{@{}llll@{}}
          \toprule
          \textbf{Magnitud} & \textbf{Unidad} & \textbf{Equivalencia}\\
          \midrule
          \multirow{3}{*}{Longitud}
            & \unidad{km}  & $1\ \text{km} = 1\,000\ \text{m}$\\
            & \unidad{m}   & $1\ \text{m} = 100\ \text{cm}$\\
            & \unidad{cm}  & $1\ \text{cm} = 10\ \text{mm}$\\[4pt]
          \multirow{3}{*}{Superficie}
            & \unidad{km\textsuperscript{2}}
                           & $1\ \text{km}^2 = 1\,000\,000\ \text{m}^2$\\
            & \unidad{ha}  & $1\ \text{ha} = 10\,000\ \text{m}^2$\\
            & \unidad{m\textsuperscript{2}}
                           & unidad base\\[4pt]
          \multirow{2}{*}{Masa}
            & \unidad{t}   & $1\ \text{t} = 1\,000\ \text{kg}$\\
            & \unidad{g}   & $1\ \text{kg} = 1\,000\ \text{g}$\\[4pt]
          \multirow{2}{*}{Tiempo}
            & \unidad{h}   & $1\ \text{h} = 3\,600\ \text{s}$\\
            & \unidad{min} & $1\ \text{min} = 60\ \text{s}$\\
          \bottomrule
        \end{tabular}}
      \end{center}
    \end{column}
    \begin{column}{0.40\textwidth}
      \begin{formulabox}[AlmAzulM]
        \AlmFontSmall
        \[
          1\ \text{ha} = 10\,000\ \text{m}^2
        \]
        \[
          1\ \text{km} = 1\,000\ \text{m}
        \]
        \[
          1\ \text{km}^2 = 100\ \text{ha}
        \]
      \end{formulabox}
      \vspace{6pt}
      \begin{tikzpicture}[scale=0.80]
        % Escalera de conversión de longitud
        \foreach \i/\lbl/\col in {%
          0/{mm}/AlmGris,
          1/{cm}/AlmAzulC,
          2/{dm}/AlmAzulC,
          3/{m}/AlmAzulM,
          4/{dam}/AlmAzulM,
          5/{hm}/AlmAzulM,
          6/{km}/AlmAzul}{
          \fill[\col, rounded corners=2pt]
            (\i*0.82, \i*0.38)
            rectangle (\i*0.82+2.6, \i*0.38+0.34);
          \node[font=\AlmFontTiny\bfseries, text=white]
            at (\i*0.82+1.3, \i*0.38+0.17) {\lbl};
        }
        \node[font=\AlmFontTiny, color=AlmGris] at (1.3,-0.28)
          {$\times 10$\quad $\times 100$\quad $\times 1000$};
      \end{tikzpicture}
    \end{column}
  \end{columns}
\end{frame}

% ---------------------------------------------------------------
\seccionframe{4 y 5. Ejemplos reales de Almería}
% ---------------------------------------------------------------

\begin{frame}{La Rambla de Almería — Longitud y anchura}
  \framesubtitle{Magnitud: Longitud · Unidades: m, km}
  \begin{columns}[c]
    \begin{column}{0.52\textwidth}
      \begin{ejemplobox}
        \AlmFontSmall
        La \textbf{Rambla de Almería} es el paseo central
        de la ciudad.\\[6pt]
        \datoalm{Longitud}{370 m}\\
        \datoalm{Anchura máxima}{50 m}\\[6pt]
        Ambas medidas expresan la magnitud
        \textbf{longitud} en metros (\unidad{m}).
      \end{ejemplobox}
      \vspace{6pt}
      \begin{pasobox}[1]
        \AlmFontFootnote
        Convertir longitud a kilómetros:
        \[
          370\ \text{m} \div 1\,000 = \mathbf{0{,}37\ \text{km}}
        \]
      \end{pasobox}
    \end{column}
    \begin{column}{0.44\textwidth}
    
          \begin{pasobox}[2]
            \AlmFontFootnote
            Estimar área (rectángulo aproximado):
            \[
              A \approx 370 \times 50 = \mathbf{18\,500\ \text{m}^2}
            \]
          \end{pasobox}
      \begin{center}
        \begin{tikzpicture}[scale=0.82]
          % Representación esquemática de la Rambla
          \fill[AlmAzulC!20, rounded corners=3pt]
            (0,0) rectangle (5.2,1.4);
          \fill[AlmVerde!35, rounded corners=2pt]
            (0.12,0.12) rectangle (5.08,1.28);
          % Acotación longitud
          \draw[acota] (0,-0.28) -- (5.2,-0.28);
          \node[font=\AlmFontTiny\bfseries, color=AlmAzulM]
            at (2.6,-0.52) {370 m};
          % Acotación anchura
          \draw[acota] (5.38,0) -- (5.38,1.4);
          \node[font=\AlmFontTiny\bfseries, color=AlmAzulM,
                rotate=90, anchor=south]
            at (5.62,0.7) {50 m};
          % Etiqueta
          \node[font=\AlmFontScript\bfseries, color=AlmAzul]
            at (2.6,0.7) {Rambla de Almería};
          % Flecha norte
          \node[font=\AlmFontTiny, color=AlmGris] at (0.35,1.15)
            {N$\uparrow$};
        \end{tikzpicture}
      \end{center}
      \vspace{4pt}
      \begin{notabox}[Tipo de información]
        \AlmFontFootnote
        \textbf{Longitud} (370 m) $\rightarrow$ \emph{Medida}\\
        \textbf{Anchura} (50 m) $\rightarrow$ \emph{Medida}\\
        \textbf{Visitantes diarios} $\rightarrow$ \emph{Dato}
      \end{notabox}
    \end{column}
  \end{columns}
\end{frame}

\begin{frame}{El Puerto de Almería — Conversión de superficie}
  \framesubtitle{Magnitud: Superficie · 1 ha = 10\,000 m²}
  \begin{columns}[c]
    \begin{column}{0.50\textwidth}
      \begin{ejemplobox}
        \AlmFontSmall
        El \textbf{Puerto de Almería} ocupa
        \textbf{35 ha} de superficie.\\[6pt]
        ¿Cuántos metros cuadrados son?\\
        ¿Cuántos km\textsuperscript{2} son?
      \end{ejemplobox}
      \vspace{6pt}
      \begin{pasobox}[1]
        \AlmFontFootnote
        Convertir a \textbf{m\textsuperscript{2}}:
        \[
          35\ \text{ha} \times 10\,000 = \mathbf{350\,000\ \text{m}^2}
        \]
      \end{pasobox}
      \vspace{4pt}
      \begin{pasobox}[2]
        \AlmFontFootnote
        Convertir a \textbf{km\textsuperscript{2}}:
        \[
          350\,000\ \text{m}^2 \div 1\,000\,000 = \mathbf{0{,}35\ \text{km}^2}
        \]
      \end{pasobox}
    \end{column}
    \begin{column}{0.46\textwidth}
      \begin{center}
        \begin{tikzpicture}[scale=0.78]
          % Cuadrado grande = Puerto (35 ha)
          \fill[AlmAzulC!30, rounded corners=4pt]
            (0,0) rectangle (4.5,4.5);
          \node[font=\AlmFontScript\bfseries, color=AlmAzul]
            at (2.25,4.1) {Puerto de Almería};
          \node[font=\AlmFontSmall\bfseries, color=AlmAzulM]
            at (2.25,2.4) {35 ha};
          \node[font=\AlmFontFootnote, color=AlmGris]
            at (2.25,1.85) {$= 350\,000\ \text{m}^2$};
          \node[font=\AlmFontFootnote, color=AlmGris]
            at (2.25,1.3)  {$= 0{,}35\ \text{km}^2$};
          % Cuadradito = campo de fútbol ~0,7 ha
          \fill[AlmVerde!50, rounded corners=2pt]
            (0.15,0.15) rectangle (1.35,1.0);
          \node[font=\AlmFontTiny, color=AlmVerde!70!black]
            at (0.75,0.57) {\(\approx\)campo};
          % Acotación
          \draw[acota] (0,-0.28) -- (4.5,-0.28);
          \node[font=\AlmFontTiny, color=AlmGris] at (2.25,-0.52)
            {escala aprox.};
        \end{tikzpicture}
      \end{center}
      \vspace{2pt}
      \begin{notabox}[Curiosidad]
        \AlmFontFootnote
        35 ha equivalen a unos
        \textbf{49 campos de fútbol}.
      \end{notabox}
    \end{column}
  \end{columns}
\end{frame}

% ---------------------------------------------------------------
\seccionframe{6. Ejercicios multinivel — Bloom}
% ---------------------------------------------------------------

\begin{frame}{Ejercicio RECORDAR}
  \bloomtag{RECORDAR}\quad\dificultad{1}
  \vspace{4pt}
  \begin{recordarbox}
    \AlmFontSmall
    \textbf{Tarea:} Completa la siguiente tabla de memoria
    (o con tus apuntes).\\[6pt]
    \begin{tabular}{@{}lll@{}}
      \toprule
      \textbf{Magnitud} & \textbf{Unidad base} & \textbf{Ejemplo en Almería}\\
      \midrule
      Longitud     & \ldots\ldots & \ldots\ldots\ldots\ldots\\[4pt]
      Superficie   & \ldots\ldots & \ldots\ldots\ldots\ldots\\[4pt]
      Masa         & \ldots\ldots & \ldots\ldots\ldots\ldots\\[4pt]
      Tiempo       & \ldots\ldots & \ldots\ldots\ldots\ldots\\[4pt]
      Temperatura  & \ldots\ldots & \ldots\ldots\ldots\ldots\\
      \bottomrule
    \end{tabular}
    \\[6pt]
    Relaciona cada magnitud con un lugar o elemento real de Almería.
  \end{recordarbox}
\end{frame}

\begin{frame}{Ejercicio COMPRENDER}
  \bloomtag{COMPRENDER}\quad\dificultad{2}
  \vspace{4pt}
  \begin{comprenderbox}
    \AlmFontSmall
    La \textbf{playa de El Zapillo} mide 2,5 km de longitud.\\[6pt]
    Responde razonando cada respuesta:
    \begin{enumerate}
      \item ¿Qué \textbf{magnitud} representa esa medida?
            ¿A qué tipo de información pertenece
            (ubicación / medida / dato)?
      \item Convierte 2,5 km a \textbf{metros}.
            Muestra la operación completa.
      \item Para calcular la \textbf{superficie} de la playa,
            ¿qué otra medida necesitarías?
            ¿Qué operación harías?
    \end{enumerate}
  \end{comprenderbox}
  \vspace{4pt}
  \begin{notabox}[Pista]
    \AlmFontFootnote
    Longitud $\times$ anchura $=$ superficie (si es rectangular).
  \end{notabox}
\end{frame}

\begin{frame}{Ejercicio APLICAR}
  \bloomtag{APLICAR}\quad\dificultad{3}
  \vspace{4pt}
  \begin{aplicarbox}
    \AlmFontFootnote
    Clasifica cada dato de la tabla como \textbf{Ubicación},
    \textbf{Medida} o \textbf{Dato}. Convierte las unidades
    cuando se indique.\\[4pt]
    \begin{tabular}{@{}p{4.2cm}lp{2.4cm}@{}}
      \toprule
      \textbf{Información} & \textbf{Tipo} & \textbf{Conversión}\\
      \midrule
      Alcazaba: 36,84°\,N / 2,47°\,O
        & \ldots & —\\[3pt]
      Catedral: 55 m de altura
        & \ldots & \ldots en km\\[3pt]
      Temperatura máx. agosto: 38 °C
        & \ldots & —\\[3pt]
      Puerto: 35 ha
        & \ldots & \ldots en m\textsuperscript{2}\\[3pt]
      Turistas 2023: 5\,800\,000
        & \ldots & —\\[3pt]
      Paseo Marítimo: 3 km
        & \ldots & \ldots en m\\
      \bottomrule
    \end{tabular}
  \end{aplicarbox}
\end{frame}

% ---------------------------------------------------------------
%  EVIDENCIA
% ---------------------------------------------------------------
\begin{frame}{Evidencia del día}
  \begin{evidenciabox}
    Entrega tu \textbf{Tabla Clasificada} con al menos
    \textbf{5 elementos reales de Almería},
    una fila por elemento:
    \vspace{6pt}
    \begin{center}
      \AlmFontSmall
      \begin{tabular}{@{}lll@{}}
        \toprule
        \textbf{Tipo (Ubicación / Medida / Dato)}
          & \textbf{Ejemplo de Almería}
          & \textbf{Unidad}\\
        \midrule
        Medida & Longitud de la Rambla: 370 m & \unidad{m}\\
        \ldots & \ldots & \ldots\\
        \bottomrule
      \end{tabular}
    \end{center}
    \vspace{6pt}
    \begin{itemize}
      \item Incluye al menos \textbf{una conversión de unidades}
            con los pasos completos.
      \item Justifica en una frase por qué cada elemento
            es Ubicación, Medida o Dato.
    \end{itemize}
  \end{evidenciabox}
\end{frame}

% ---------------------------------------------------------------
%  ERROR FRECUENTE
% ---------------------------------------------------------------
\begin{frame}{¡Cuidado con este error!}
  \begin{errorbox}
    \begin{columns}[c]
      \begin{column}{0.50\textwidth}
        \incorrecto\ \textbf{Error frecuente}\\[4pt]
        Confundir un \textbf{dato} (varía con el tiempo)
        con una \textbf{medida} (valor fijo).\\[6pt]
        «La temperatura de Almería es
        \cancel{una medida fija}» — \textbf{incorrecto}.\\
        La temperatura cambia: es un \textbf{dato}.
      \end{column}
      \begin{column}{0.46\textwidth}
        \correcto\ \textbf{Correcto}\\[4pt]
        La \textbf{altura} de la Catedral (55 m)
        \emph{no cambia con el tiempo}:
        es una \textbf{medida}.\\[6pt]
        La \textbf{temperatura} máxima de un día
        \emph{sí cambia} cada día:
        es un \textbf{dato}.
      \end{column}
    \end{columns}
  \end{errorbox}
  \vspace{8pt}
  \begin{notabox}[Regla rápida]
    \AlmFontSmall
    ¿Cambiaría si volvemos a medirlo mañana?
    \correcto\;Sí $\Rightarrow$ \textbf{dato}.\quad
    \correcto\;No $\Rightarrow$ \textbf{medida}.
  \end{notabox}
\end{frame}

% ---------------------------------------------------------------
%  RESUMEN
% ---------------------------------------------------------------
\begin{frame}{Resumen de la sesión 02}
  \begin{resumenbox}
    \begin{itemize}
      \item La información matemática de una ciudad se clasifica
            en \textbf{Ubicación}, \textbf{Medida} y \textbf{Dato}.
      \item Toda medida tiene \textbf{valor numérico} y \textbf{unidad};
            escribir la unidad es obligatorio.
      \item Las magnitudes más frecuentes son longitud (\unidad{m}, \unidad{km}),
            superficie (\unidad{m\textsuperscript{2}}, \unidad{ha}),
            masa (\unidad{kg}), tiempo (\unidad{s}, \unidad{h}) y
            temperatura (\unidad{°C}).
      \item Conversiones clave:
            $1\ \text{km}=1\,000\ \text{m}$;
            $1\ \text{ha}=10\,000\ \text{m}^2$.
      \item Un dato \textbf{cambia} con el tiempo;
            una medida \textbf{no cambia}.
    \end{itemize}
  \end{resumenbox}
  \vspace{6pt}
  \begin{center}
    {\AlmFontSmall\color{AlmAzulM}
    \faArrowCircleRight\enskip
    \textbf{Próxima sesión:}
    Plano por cuadrantes — localizar lugares con letra y número.}
  \end{center}
\end{frame}

\end{document}



